%% LyX 2.3.4.2 created this file.  For more info, see http://www.lyx.org/.
%% Do not edit unless you really know what you are doing.
\documentclass[11pt,spanish]{article}
\renewcommand{\rmdefault}{cmr}
\usepackage[T1]{fontenc}
\usepackage[utf8]{inputenc}
\usepackage{geometry}
\geometry{verbose,tmargin=2cm,bmargin=2cm,lmargin=2cm,rmargin=2cm}
\setlength{\parskip}{\bigskipamount}
\setlength{\parindent}{0pt}
\usepackage{amstext}
\usepackage{amsthm}
\usepackage{amssymb}

\makeatletter
%%%%%%%%%%%%%%%%%%%%%%%%%%%%%% Textclass specific LaTeX commands.
\theoremstyle{definition}
\newtheorem{defn}{\protect\definitionname}[section]
\theoremstyle{definition}
\newtheorem{xca}{\protect\exercisename}[section]

%%%%%%%%%%%%%%%%%%%%%%%%%%%%%% User specified LaTeX commands.
\usepackage{titling}
\setlength{\droptitle}{-2cm}
\usepackage{multicol}
\usepackage{enumitem}
\usepackage{proof}
\usepackage{mathrsfs}
\usepackage{amsfonts}
\usepackage{forest}
\usepackage{ebproof}
\usepackage{amssymb}
\definecolor{Guinda}{rgb}{0.49, 0.05, 0.12}
\definecolor{Rojo}{rgb}{0.8, 0, 0}
\definecolor{Verde}{rgb}{0,0.4,0}
\definecolor{Azul}{rgb}{0,0.4,0.4}
\definecolor{Naranja}{rgb}{0.8,0.4,0}

\makeatother

\usepackage{babel}
\addto\shorthandsspanish{\spanishdeactivate{~<>.}}

\usepackage{listings}
\lstset{mathescape=true}
\providecommand{\definitionname}{Definición}
\providecommand{\exercisename}{Ejercicio}

\begin{document}
\title{Teoría de la Computación, 2021-1\\
Nota de clase 1: Introducción al curso\thanks{Material mayormente tomado de \cite{key-1} con ligeras modificaciones.}}
\author{Manuel Soto Romero}
\date{\today \\ Colegio de Ciencia y Tecnología UACM}

\maketitle
\renewcommand*{\thefootnote}{\arabic{footnote}}
\setcounter{footnote}{0}
\setcounter{section}{1}
\renewcommand*\lstlistingname{C\'odigo}

El curso de Teoría de la Computación tiene como fin que apliques las
teorías de lenguajes formales, gramáticas y las máquinas de estados
finitos, en el diseño e implementación de aplicaciones reales de software.
Además que seas capaz reconocer aquellos problemas que no puedan ser
resueltos utilizando una computadora. 

En esta nota trataremos de entender por qué todo esto es importante
en el perfil de cualquier carrera de computación y daremos una breve
introducción a algunos de los elementos básicos que manipularemos
a lo largo del curso. Estudiaremos con mayor detalle estos conceptos
en otras notas.

\subsection{¿Por qué estudiar Teoría de la Computación?}

Quizá una forma de entender por qué es importante para nosotros el
estudio de esta teoría sea platicando algunas de las aplicaciones
acerca de los temas que estudiaremos a lo largo de estas notas. \cite{key-2}
\begin{itemize}
\item \textbf{Autómatas Finitos. }Constituyen un modelo útil para muchos
tipos de hardware y software. Por mencionar algunos de estos modelos
tenemos:
\begin{enumerate}
\item Software para diseñar y probar el comportamiento de circuitos digitales.
\item El analizador léxico de un compilador.
\item Software para explorar patrones en cadenas.
\item Software para verificar sistemas de todo tipo y que tengan un número
finito de estados diferentes, por ejemplo, protocolos de intercambio
de información.
\end{enumerate}
\item \textbf{Gramáticas. }Son modelos útiles en el diseño de software que
sirve para procesar datos con una estructura recursiva. El ejemplo
más conocido de éstas se puede encontrar en los Analizadores Sintácticos,
los cuales son componentes de un compilador (o intérprete) encargados
de verificar que lo que escribimos en un determinado lenguaje de programación
está bien escrito. Por ejemplo, una regla gramatical E := E + E nos
dice que una expresión puede formarse tomando dos expresiones y conectándolas
poniendo un símbolo de suma entre ellos.
\item \textbf{Expresiones Regulares. }También nos permiten revisar la estructura
de los datos, especialmente de las cadenas. Los patrones de una cadena
que puede describir una expresión regular, también pueden ser descritos
por los autómatas finitos. Por ejemplo, los sistemas operativos basados
en UNIX, permiten definir expresiones regulares, como la siguiente:
\begin{center}
\texttt{'{[}A-Z{]} {[}a-z{]}{*} {[}{]} {[}A-Z{]}{[}A-Z{]}'}
\par\end{center}

esta expresión regular representa palabras que comienzan por una letra
mayúscula seguida de un espacio y dos letras mayúsculas. Por ejemplo
\texttt{``México DF''}.
\item \textbf{Decidibilidad. }Nos permite responder la pregunta ¿Qué puede
hacer una computadora?, los problemas que una computadora puede resolver,
son llamados decidibles.
\item \textbf{Intratabilidad. }Nos permite responder la pregunta ¿Qué puede
hacer una computadora de manera eficiente? Estudia los problemas que
una computadora puede resolver en un tiempo ``descente'' o mejor
dicho ``tratable''. Esto también se estudia con detalle en los cursos
de Análisis de Algoritmos.
\end{itemize}

\subsection{Conceptos básicos}

Se presentan a continuación las definiciones de los términos más importantes
empleados en la Teoría de la Computación. En general, manejaremos
tres conceptos:
\begin{itemize}
\item Alfabetos: Conjunto de símbolos
\item Cadenas: Lista de símbolos de un alfabeto
\item Lenguaje: Conjunto de cadenas de un mismo alfabeto.
\end{itemize}
Definamos esto con más formalidad.

\subsubsection{Alfabetos y cadenas}

\bigskip{}

\begin{defn}
\emph{(Alfabeto) Un alfabeto es un conjunto de símbolos finito y no
vacío. Denotamos a un alfabeto con la letra griega $\Sigma$. }
\end{defn}
%
Algunos ejemplos de alfabetos son:
\begin{itemize}
\item $\Sigma=\left\{ 0,1\right\} $, el alfabeto binario.
\item $\Sigma=\left\{ a,b,...,z\right\} $, el conjunto de todas las letras
minúsculas.
\item El conjunto de todos los caracteres ASCII.
\end{itemize}
\bigskip{}

\begin{defn}
\emph{(Cadena) Una cadena es una secuencia finita de símbolos seleccionados
de algún alfabeto.} 
\end{defn}
%
Por ejemplo, 01101 y 111 son cadenas del alfabeto binario.

\bigskip{}

\begin{xca}
Dado el siguiente alfabeto $\Sigma=\left\{ 1,2,3,4,5,6,a,b,c,d\right\} ,$da
cinco ejemplos de cadenas que puedas formar con él.
\end{xca}
%
\bigskip{}

\begin{defn}
\emph{(Cadena vacía) La cadena vacía es aquella que no tiene ningún
símbolo, es denotada por $\varepsilon$ y está presente en todo alfabeto.}
\end{defn}
%
\bigskip{}

\begin{defn}
\emph{(Longitud de una cadena) La longitud de una cadena $w$ es el
número de símbolos que la conforman, se denota por $|w|$.}
\end{defn}
%
Ejemplos:
\begin{itemize}
\item $|01101|=5$
\item $|011|=3$
\item $|\varepsilon|=0$.
\end{itemize}
\bigskip{}

\begin{xca}
Calcula la longitud de todas las cadenas que definiste en el Ejercicio
1.1.
\end{xca}
%
\bigskip{}

\begin{defn}
\emph{(Reversa de una cadena) La reversa de una cadena $w$ es simplemente
la representación en orden inverso de los símbolos que la conforman,
se denota por $w^{R}$.}
\end{defn}
%
Ejemplos:
\begin{itemize}
\item \emph{$0101^{R}=1010$}
\item $abc^{R}=cba$
\item $\varepsilon^{R}=\varepsilon$
\end{itemize}
\bigskip{}

\begin{xca}
Calcula la reversa de todas las cadenas que definiste en el Ejercicio
1.1.
\end{xca}
%
\bigskip{}

\begin{defn}
\emph{(Potencia de un alfabeto) La potencia $k$ de un alfabeto ($\Sigma^{k}$)
se define como el conjunto de las cadenas de longitud $k$, tales
que cada uno de los símbolos de las mismas pertenecen a $\Sigma$.} 
\end{defn}
%
Por ejemplo, si tomamos $\Sigma=\left\{ 0,1\right\} $, tenemos las
siguientes potencias:
\begin{itemize}
\item $\Sigma^{0}=\left\{ \varepsilon\right\} $
\item $\Sigma^{1}=\left\{ 0,1\right\} $
\item $\Sigma^{2}=\left\{ 00,01,10,11\right\} $
\item $\Sigma^{3}=\left\{ 000,001,010,011,100,101,110,111\right\} $
\item ...
\end{itemize}
Para denotar el conjunto de todas las cadenas de un alfabeto, usamos
la notación \emph{estrella, }también llamada \emph{Cerradura de Kleen}.\emph{
}Por ejemplo, para nuestro alfabeto binario, tenemos:

\[
\Sigma^{*}=\left\{ 0,1\right\} ^{*}=\left\{ \varepsilon,0,1,00,01,10,11,000,...\right\} 
\]

dicho de otro modo, es el conjunto que tiene a las cadenas de longitud
0, a las de longitud 1, a las de 2 y así sucesivamente.

\[
\Sigma^{*}=\Sigma^{0}\cup\Sigma^{1}\cup\Sigma^{2}\cup...
\]

En ocasiones necesitaremos descartar a $\varepsilon$ de nuestro conjunto
de cadenas, esto es denotado por $\Sigma^{+}$, y es llamada \emph{Cerradura
Positiva} es decir:

\[
\Sigma^{+}=\Sigma^{1}\cup\Sigma^{2}\cup...
\]

\bigskip{}

\begin{xca}
Dado el siguiente alfabeto $\Sigma=\left\{ a,b,c\right\} $, obtén
$\Sigma^{0},\Sigma^{1},\Sigma^{2}$ y muestra los primeros 16 elementos
de $\Sigma^{*}$y de $\Sigma^{+}$.
\end{xca}
%
\bigskip{}

\begin{defn}
\emph{(Concatenación de cadenas) Definimos la concatenación de dos
cadenas $x$ y $y$ como la cadena formada por una copia de $x$ seguida
de una copia de $y$. Lo denotamos poniendo las cadenas juntas, es
decir $xy$. Es decir, si $x=a_{0}a_{1}a_{2}...a_{i}$ y $y=b_{0}b_{1}b_{2}...b_{j}$
entonces $xy=a_{0}a_{1}a_{2}...a_{i}b_{0}b_{1}b_{2}...b_{j}$. Notemos
además que $|xy|=i+j$.} 
\end{defn}
%
Por ejemplo, si $x=01101$ y $y=110$, entonces $xy=01101110$ y $yx=11001101$.

Es importante destacar además que $\varepsilon$ es el neutro de la
concatenación, es decir, si lo concatenamos con cualquier cadena $s$
siempre obtendremos $s$.

\[
\varepsilon s=s\varepsilon=s
\]


\subsubsection{Lenguajes}

Ahora que sabemos qué es un alfabeto y qué significa que una cadena
pertenezca a él. Veamos lo que es un \emph{lenguaje}. 

Se le denomina lenguaje a un subconjunto $L$ de cadenas tomadas de
$\Sigma^{*}$ ($L\subseteq\Sigma^{*})$. Es decir, dadas todas las
cadenas que podemos formar con un alfabeto, tomamos algunas de ellas
para conformar un lenguaje. Veamos algunos ejemplos:
\begin{itemize}
\item El español es un lenguaje formado por todas las cadenas escritas en
español tomadas del abecedario. 
\item El lenguaje de programación \textsc{Haskell} es un lenguaje formado
por todos los símbolos y palabras reservadas necesarias para escribir
un programa.
\item El lenguaje de todas las cadenas que constan de $n$ ceros, seguidos
de $n$ unos para cualquier $n\geq0:$

\[
\left\{ \varepsilon,01,0011,000111,...\right\} 
\]

\item El conjunto de cadenas formadas por el mismo número de ceros que de
unos:

\[
\left\{ \varepsilon,01,10,0011,0101,1001,...\right\} 
\]

\item El conjunto de números binarios cuyo valor es un número primo:

\[
\left\{ 10,11,101,111,1011,...\right\} 
\]

\item $\Sigma^{*}$ es un lenguaje para cualquier alfabeto $\Sigma$.
\item $\varnothing$, el lenguaje vacío, es un lenguaje de cualquier alfabeto.
\item $\left\{ \varepsilon\right\} $, el lenguaje que consta sólo de la
cadena vacía, también es un lenguaje de cualquier alfabeto. (Observa
que $\varnothing\neq\varepsilon$).
\end{itemize}
Dado que un lenguaje es un conjunto, podemos describir lenguajes usando
la notación por comprensión de los conjuntos:

\[
\left\{ s:\text{algo acerca de }s\right\} 
\]

Por ejemplo:
\begin{itemize}
\item $\left\{ s:s\text{ consta de un número igual de ceros que de unos}\right\} $
\item $\left\{ s:s\text{ es un entero binario que es primo}\right\} $
\item $\left\{ s:s\text{ es un programa en \textsc{Haskell} sintácticamente correcto}\right\} $
\end{itemize}
También es común que las definiciones dependan de ciertos parámetros
y condiciones sobre estos. Por ejemplo:
\begin{itemize}
\item $\left\{ 0^{n}1^{n}:n\geq1\right\} $
\item $\left\{ 0^{i}0^{j}:0\leq i\leq j\right\} $
\end{itemize}
Dado que los lenguajes son conjuntos, todas las operaciones de conjuntos
aplican, tales como $L\cup M$, $L\cap M,L-M,L^{c}=\Sigma^{*}-L$.
Adicionalmente tenemos otras operaciones.

\bigskip{}

\begin{xca}
Sea $\Sigma=\left\{ a,b,c\right\} $ y $L_{1}=\left\{ \varepsilon,a,ab\right\} $
y $L_{2}=\left\{ b,ac,ab\right\} $. Obtén:
\end{xca}
\begin{itemize}
\item $L_{1}\cup L_{2}$
\item $L_{1}\cap L_{2}$
\item $L_{1}-L_{2}$
\end{itemize}
%
\bigskip{}

\begin{defn}
\emph{(Concatenación de lenguajes) La concatenación de lenguajes se
define como}

\[
LM=\left\{ uv:u\in L\land v\in M\right\} 
\]
\end{defn}
%
Por ejemplo, si $A=\{hola,adi\acute{o}s\}$ y $B=\left\{ casa\right\} $,
entonces $AB=\left\{ holacasa,adioscasa\right\} $

\bigskip{}

\begin{xca}
Usando el alfabeto y lenguajes del Ejercicio 1.5, obtén $L_{1}L_{2}$.
\end{xca}
%
\bigskip{}

\begin{defn}
\emph{(Reversa de un lenguaje) La reversa de un lenguaje se define
como}

\[
L^{R}=\left\{ w^{R}:w\in L\right\} 
\]
\end{defn}
%
Por ejemplo, si $L=\{01,10,1011\}$, entonces $L^{R}=\left\{ 10,01,1101\right\} $

\bigskip{}

\begin{xca}
Usando el alfabeto y lenguajes del Ejercicio 1.5, obtén $L_{1}^{R}$
y $L_{2}^{R}$
\end{xca}
%
\bigskip{}

\begin{defn}
\emph{(Potencia de un lenguaje) La potencia $k$ de un lenguaje ($L^{k}$)
se define como el conjunto que contiene la concatenación de L consigo
mismo $k$ veces}.
\end{defn}
%
Algunos ejemplos sobre el lenguaje $A=\left\{ ab\right\} $
\begin{itemize}
\item $A^{0}=\left\{ \varepsilon\right\} $
\item $A^{1}=A=\left\{ ab\right\} $
\item $A$$^{2}=AA=\left\{ abab\right\} $
\item $A^{3}=AAA=\left\{ ababab\right\} $
\item ...
\end{itemize}
Y por supuesto, también tiene sus respectivas cerraduras de Kleen
y Positiva.

\[
L^{*}=L^{0}\cup L^{1}\cup L^{3}\cup...
\]

\[
L^{+}=L^{1}\cup L^{2}\cup...
\]

\bigskip{}

\begin{xca}
Usando el alfabeto y lenguajes del Ejercicio 1.5, obtén los primeros
11 elementos de $L_{1}^{*}$.
\end{xca}
%
En la Teoría de la Computación, un \emph{problema} es la cuestión
de decidir si una determinada cadena pertenece a un lenguaje. Como
estudiaremos más adelante en otra nota, cualquier cosa que denominemos
problema, puede ser expresado como elemento de un lenguaje. Dicho
de otra forma, si $\Sigma$ es un alfabeto y $L$ es un lenguaje de
$\Sigma$, entonces el problema $L$ es:
\begin{center}
\emph{Dada una cadena $s$ de $\Sigma^{*}$, decidir si $s$ pertenece
o no a $L$}
\par\end{center}
\begin{thebibliography}{1}
\bibitem{key-1} Hopcroftt, J. E; Montwani, R., and Ullman, J. \emph{Introduction
to Automata Theory Language, and Computation}. Addison Wesley, third
edition 2007. 
\end{thebibliography}

\end{document}
